\documentclass[a6paper,fontsize=10pt,twoside,openright]{scrbook}
\usepackage[a6paper,top=1cm,bottom=1.9cm,left=1cm,right=1cm,footskip=1cm,headsep=0.2cm]{geometry}
\usepackage[swedish]{babel}
\usepackage{graphicx,tocloft,metalogo,titletoc,fontspec,titlesec,eso-pic,scrpage2,anyfontsize,microtype,array,marvosym,mdframed,textcomp,amssymb}
\usepackage[columns=1]{idxlayout}

\usepackage[htt]{hyphenat}

\usepackage{setspace}
\renewcommand{\baselinestretch}{0.98} 

%\usepackage{showframe}
\usepackage{makeidx}
\makeindex

\let\oldemptyset\emptyset
\let\emptyset\varnothing

%% eso-pic for background on chapter pages %%
\newcommand\BackgroundPic{\includegraphics[width=\paperwidth,height=\paperheight,keepaspectratio]{elements/chapter2.pdf}}

%% scrpage2 pagestyling %%
\pagestyle{scrheadings}
\clearscrheadfoot

%% change empty page produced by \cleardoublepage into scrheading style %%
\makeatletter
\let\ps@empty\ps@scrheadings
\makeatother

\newcommand{\rpt}{
        \raisebox{.2ex}{:}
        \hspace*{-0.5ex}\raisebox{-.4ex}{\rule{.1ex}{2.5ex}\,\rule{.1ex}{2.5ex}}}

\newcommand{\revrpt}{
        \raisebox{-.4ex}{\rule{.1ex}{2.5ex}\,\rule{.1ex}{2.5ex}}
        \hspace*{-0.5ex}\raisebox{.2ex}{: }}

%% chapter and section styling %%
\titleformat{\chapter}[hang]{}{}{0pt}{\centering\fontsize{20}{70}\selectfont\chapterfont\texttt}
\titleformat{\section}[hang]{}{}{0pt}{\textbf}
\titlespacing*{\section}{0pt}{0pt}{0pt}
\titlecontents{chapter}[0em]{\addvspace{3pt}}{\contentslabel{}}{}
              {\titlerule*[0.4pc]{.}\contentspage}
\newfontfamily\chapterfont[]{ChopinScript}

%% Command for chapter page, will remove foot and head %%
\newcommand{\chapterpage}[2]{
  \newpage
  \null
  \cleardoublepage
  %\AddToShipoutPicture*{\BackgroundPic}
  \chapter[#1]{#2}
  \renewcommand{\leftmark}{#1}
  \newpage
  \ohead{\textsc{\scriptsize\leftmark}}
  \ofoot[\pagemark]{\scriptsize\pagemark}
  %\cfoot{\centering\includegraphics[keepaspectratio]{elements/footelement2.pdf}}
}
\newcommand{\hiddenchapterpage}[1]{
  \newpage
  \null
  \cleardoublepage
  %\AddToShipoutPicture*{\BackgroundPic}
  \chapter*{#1}
  \renewcommand{\leftmark}{#1}
  \newpage
  \ohead{\textsc{\scriptsize\leftmark}}
  \ofoot[\pagemark]{\scriptsize\pagemark}
  %\cfoot{\centering\includegraphics[keepaspectratio]{elements/footelement2.pdf}}
}
\makeatletter
\renewcommand\chapter{\if@openright\cleardoublepage\else\clearpage\fi
  \thispagestyle{empty}
  \global\@topnum\z@
  \@afterindentfalse
  \ofoot[\pagemark]{\scriptsize\pagemark}
  \cfoot{}
  \ohead{}
  \secdef\@chapter\@schapter}
\makeatother

%% Command for songs %%
\newcommand{\visa}[1]{
  \index{#1}
  \section{#1}
}
\newcommand{\fvisa}[2]{
  \index{#1}\index{#2}
  \section{#1}
}
\newcommand{\svisa}[3]{
  \index{#1}\index{#2}
  \section{#3}
}

%% TOC %%
\setcounter{tocdepth}{0}
\tocloftpagestyle{scrheadings}

%% MISC %%
\pagenumbering{arabic}
\renewcommand{\chaptername}{}
\renewcommand{\thechapter}{}
\renewcommand{\thesection}{}

\makeatletter
\renewenvironment{theindex}
                 {{\noindent\small\textrm{\textbf{TITEL och FÖRSTARADEN REGISTER}}\vspace*{10pt}\newline\par\textbf{\normalsize Symboler}\par}
                   \@mkboth{\MakeUppercase\indexname}
                           {\MakeUppercase\indexname}
                           \let\item\@idxitem}
                 {}
\makeatother

\newmdenv[
  topline=false,
  bottomline=false,
  rightline=false,
  linewidth=1pt,
  skipabove=5pt,
  skipbelow=3pt,
  innertopmargin=0pt,
  innerbottommargin=0pt
]{leftborder}
\makeatletter

\newcommand*\cleartoleftpage{%
  \clearpage
  \ifodd\value{page}\hbox{}\newpage\fi
}

\begin{document}
\setlength{\parindent}{15pt}
\vspace*{6.5cm}
\hspace*{1cm}
\includegraphics[keepaspectratio,width=7cm]{elements/logo.pdf}
\clearpage
\noindent Denna förträffliga oersättliga Manual tillhör:
\ohead{\textsc{\scriptsize\leftmark}}
\ofoot[\pagemark]{\scriptsize\pagemark}
\clearpage
\null
\clearpage
\newgeometry{top=1cm,bottom=1.2cm,left=1cm,right=1cm,footskip=0pt,headsep=0.2cm}
\null
\vfill
    {\noindent\small\centering
      Redaktör: ?\\
      Grafisk formgivare: ?\\
      Logotyp: Jens Bohlin\\
      Producerad med \XeTeX\\
      v1.0: 1994-11-21 (1000 exemplar)\\
      v1.1: 1996-01-02 (1000 exemplar)\\
      v2.001: 2001-11-15 (1000 exemplar)\\
      v11111011111: 2015-??-?? (1200 exemplar)\\
      Tryckeri: Lyxo AB\par
    }
\restoregeometry
\cleardoublepage
%\cfoot{\centering\includegraphics[keepaspectratio]{elements/footelement2.pdf}}
\noindent{\large\textrm{\textbf{FÖRORD}}\vspace*{10pt}}\\
\indent ost
\cleardoublepage
\noindent{\large\textrm{\textbf{FÖRORD från v2.001}}\vspace*{10pt}}\\
\indent ost
\cleardoublepage
\noindent{\large\textrm{\textbf{FÖRORD från v1.0}}\vspace*{10pt}}\\
\indent Detta är en gemensam sångbok för samtliga datavetare i
Sverige; C i Lindköping, DVP i Uppsala och C i Umeå. Vi har genom
denna bok fått ytterligare en sak gemensamt.\\ \indent Att skapa en
gemensam sångbok har diskuterats flerga gånger, jag minns att
kontakter tog under 100-köret 1990 eller var det 1991?  Saker som
händer under 100-kör är inget man minns, utan saker som man senare
tror sig har hänt?! Den gången brakade arbetet samman efter några
veckor. Men idén levde kvar och gav 3 (eller 4) år senare detta som
resultat.\\ \indent Sångboken bygger på ``DVL:s Lilla blå
v3.0$\beta$'' som aldrig han bli tryckt, innan arbetet kom igång för
att skapa denna bok. Dock skall sägas att denna bok är betydligt
större än ``Lilla blå v3.0$\beta$'', omkring 200 sidor för att inte
vara exakt.\\ \indent När \textit{manualen} trycks i december 1994 är
boken den största sångbok som finns runt om i vårt rike! Ok Psalmboken
är större, men inte lika roligt att ha på gasquer!\\ \indent Vi som
varit ansvariga för denna bok, har försökt att i så stor utsträckning
som möjligt att få med melodier, låtarnas ursprung, och hur man skall
sjunga. Allt för att underlätta för folk som inte är i kontakt med
lika mycket sånger som vissa av oss andra.\\ \indent Det är många som
har hjälpt till med denna bok, inte enbart datavetare utan även
personer från helt andra linjer och orter, jag vill rikta ett
speciellt tack till Er.\\ \indent Manualens generalgenomsjungning ägde
rum på\\ Strålsnäs herrgård den 18 November 1994, då vi avnjöt en
underbar rådjursmiddag. På denna sittning sjöng vi så mycket vi orkade
och hade efter ca 4 timmar traverserat ca 40 sånger.\\ \indent Med på
denna sångboksgenomsjungning och därav lika (o)ansvariga som jag till
de eventuella felen var,\\ \textit{Micke Liljeblad, Mats Hägglund, Jakob
  Engblom, Pär\\ Mattsson, Martin Davidson, Anders Jackson, Björn\\
  Knutsson, Måns Engmem, Måns af Klercker, Jonas\\ Högström, Calle
  Englund, Andreas Gref} och sist jag själv.\\ \indent Innan jag
kommer till alla som skall tackas vill jag bara nämna att en sångbok
är till för att sjungas ur, detta betyder inte att man behöver
sjunga. Följande personer är jag evigt(?) tacksam till;\\

\noindent
\begin{tabular}{@{}p{0.5\textwidth}p{0.5\textwidth}@{}}
  Anders Jansson & Anders Wilhelm\\
  Bertil Spolander & Björn Knutsson\\
  Björn Torkelsson & Erik Åhlin\\
  Göran Svensson & Håkan Lindholm\\
  Ingela Ljungkvist & Jan-Olof Flink\\
  Jens Bohlin & Joakim Thörnkvist\\
  Jonas Högström & Lars Hägglund\\
  Måns af Klercker & Olof Appelblad\\
  Pär Mattson
\end{tabular}
\newpage
\indent Dessutom vill jag tacka alla som lusläst boken, sammt alla Ni
som gjorde förhandsbeställningar. Det finns säkert många andra som jag
missat och jag är precis lika tacksam till er som alla
andra!\\ \indent Om ni har några bidrag eller saker som kan vara värda
att rätta till så skicka dessa till någon av festerierna nedan;\\

\noindent
\begin{tabular}{@{}p{0.3\textwidth}p{0.7\textwidth}@{}}
  Piraya klubben & piraya@ts.umu.se\\
  Klubbverket & dvl-klubbverket@student.docs.uu.se\\
  CC & cc@cyd.liu.se
\end{tabular}\\

\indent Till sist, kom ihåg att hela sångboken och fellista finns på:
\texttt{http://www.ts.umu.se/\raisebox{.4ex}\texttildelow manualen/}\\ \\ Umeå den 21 November
1994\\ \\ NeNNe
\cleartoleftpage
\noindent{\large\textrm{\textbf{TRAVERSERING}}\vspace*{10pt}}\\ \indent
Genom hela denna bok gäller följande traverseringsregler. I sånger där
en försångare eller dylikt skall sjunga ensam används fet text. Ifall
något speciellt såsom en handling inträffar vid ett ord, så används
(kursiv), där ordet för handling eller beteende står inom ( ). En
sådan handling gäller oftast maximalt raden ut. Ifall en längre
förklaring behövs kursiveras texten. Beskrivningen finns sedan efter
sångens slut. I denna beskrivning kan mycket annat matnyttigt finnas,
läs den gärna.\\ \indent I vissa sånger skall könen sjunga olika
rader, i dessa sånger används {\Large\Male} för manligt kön och där
kvinnligt kön skall sjunga används givetvis då {\Large\Female}.
\begin{leftborder}
  Då det ibland kan förekomma lokala variantioner på vissa sånger, så
  används en  vertikal linje i högra kanten för att markera att detta
  är olika alternativ. Ibland betyder linjen att det är raderna som är
  alternativen, och ibland betyder linjen att det är olika
  versalternativ. Men den kloke läsaren förstår säkert när vad är vad.
\end{leftborder}
\hspace{\indent} För att göra allt ännu svårare så skall man oftast sjunga med stark
stämma där VERSALER används, men observera att man inte alltid skall
göra detta. Vem har sagt att det skall vara lätt?
\newpage
%% begin toc on odd page and style the title %%
\cleardoublepage
\renewcommand{\contentsname}{\vspace{-2.1cm}\large\textrm{INNEHÅLL}\vspace{-1.2cm}}
\tableofcontents

%% Remove the toc title from heading %%
\renewcommand{\leftmark}{}

\setlength{\parindent}{0pt}
\chapterpage{Av tradition}{Av tradition}
\visa{DU GAMLA, DU FRIA}
Du gamla, Du fria, Du fjällhöga nord\\
Du tysta, Du glädjerika sköna!\\
Jag hälsar Dig, vänaste land uppå jord,\\
||: Din sol, Din himmel, Dina ängder gröna. :||\\
\\
Du tronar på minnen från fornstora dar,\\
då ärat Ditt namn flög över jorden.\\
Jag vet att Du är och Du blir vad du var.\\
||: Ja, jag vill leva jag vill dö i Norden. :||\\
\\
Jag städs vill dig tjäna mitt älskade land,\\
din trohet till döden vill jag svära.\\
Din rätt, skall jag värna, med håg och med hand,\\
||: din fana, högt den bragderika bära. :||\\
\\
Med Gud skall jag kämpa, för hem och för härd,\\
för Sverige, den kära fosterjorden.\\
Jag byter Dig ej, mot allt i en värld\\
||: Nej, jag vill leva jag vill dö i Norden. :||\\
\\
{\footnotesize\textit{Text: Richard Dybeck, 1844\\
\\
Sveriges nationalsång av tradition sedan 1866.\\
De två sista verserna sjungs sällan.}}
\clearpage

%\fvisa{SVERIGES FLAGGA}{Flamma stolt mot dunkla skyar}
\vspace{10pt}
Flamma stolt mot dunkla skyar\\
lik en glimt av sommarens sol!\\
Över Sveriges skogar, berg och byar,\\
över vatten och viol!\\
Du som sjunger, när Du bredes\\
som vår gamla lyckas tolk.\\
Solen lyser! Solen lyser!\\
Ingen vredes åska slog vårt tappra folk!\par
\vspace{10pt}
Flamma högt vårt kärlekstecken!\\
Värm oss, när det blåser kallt!\\
Stråla ut de blåa vecken\\
kärlek mera stark än allt!\\
Sveriges flagga! Sveriges ära!\\
Fornklenod och framtidstolk!\\
Gud är med oss! Gud är med oss!\\
Han skall bära stark vårt fria svenska folk\par
\vspace{10pt}
{\footnotesize\textit{Text: K.G. Ossiannilsson\\
Musik: Hugo Alfvén}}

%\fvisa{JÄMTLANDSSÅNGEN}{Så tåga vi tillsammans bort}
\vspace{10pt}
Så tåga vi tillsammans bort\\
mellan Jämtlands gröna ängar\\
bort mellan nyland som prunka\\
fulla av bröllopsblomsters prakt.\\
Så skådom nu med gamman hän\\
emot berg i blåa fjärran hän,\\
över sjöar, strömmar, skogar\\
jämt kring bygder på vakt.\\
\\
\revrpt Fagert är landet\\
som blev vår lott och arvedel.\\
Så firom dess fägring nu\\
med sång och stråkars spel.\\
Så tändom ånyo det hopp\\
som våra fäder närt.\\
För slit och mödor\\
av fröjd och sol\\
ett mått oss beskärt\rpt\\
\\
{\footnotesize\textit{Text \& Musik: Wilhelm Peterson-Berger}}
\clearpage

%\input{songs/jamtlandssången.tex}
\input{songs/studentsången.tex}
\visa{O, GAMLA KLANG- OCH JUBELTID}
{\footnotesize\textit{Melodi: Oh alte Burschenherrlichkeit}}\par
\vspace{10pt}
O gamla klang och jubeltid\\
ditt minne skall förbliva\\
och än åt livets bistra strid,\\
ett rosigt skimmer giva.\\
Snart tystnar allt vårt yra skämt,\\
vår sång blir stum, vårt glam förstämt.\\
O, jerum, jerum, jerum.\\
O, quae mutatio rerum!\par
\vspace{10pt}
Var äro de som kunde allt,\\
blott ej sin ära svika.\\
Som voro män av äkta halt\\
och världens herrar lika?\\
De drogo bort från vin och sång\\
till vardagslivets tråk och tvång.\\
O, jerum, jerum, jerum.\\
O, quae mutatio rerum!\\
\begin{leftborder}
\begin{tabular}{l l}
  (Filosofer)  & Den ene vetenskap och vett\\		
               & in i scholares mänger,\\
  (Jurister)   & Den andre i sitt anlets svett\\
               & på paragrafer vränger,\\
  (Teologer)   & en plåstrar själen, som är skral,\\
  (Medicinare) & en lappar hop dess trasiga fodral;\\
  (Alla)       & O, jerum, jerum, jerum,\\
               & O, quae mutatio rerum!
\end{tabular}
\end{leftborder}
\newpage
Men hjärtat i en sann student,\\
kan ingen tid förfrysa.\\
Den glädjeeld, som där har tänt,\\
hans hela liv skall lysa.\\
Det gamla skalet brustit har\\
men \textit{kärnan} finnes frisk dock kvar,\\
och vad han än må mista,\\
den skall dock aldrig brista!\par
\vspace{10pt}
Så sluten, bröder, fast vår krets,\\
till glädjens värn och ära!\\
Trots allt vi tryggt och väl tillfreds,\\
vår vänskap trohet svära.\\
Lyft bägar'n högt, och klinga vän!\\
De gamla gudar leva än\\
\revrpt bland skålar och pokaler\rpt\par
\vspace{10pt}
{\footnotesize\textit{Text: August Lindh, 1925\\ Musik:
Eugen Höfling, 1825}}\par
\vspace{10pt}
{\footnotesize\textit{Vid ``kärnan' dunkas näven i bordet en
gång. Den sista versen sjungs stående samt stående på stolen, efter
sista versen sjungits sätter sig ingen åter till bords.}}

\visa{UNDER SVEA BANÈR}
\vspace{10pt}
Under Svea Banér Himlen seger oss ger;\\
Då för Konung och Land Äran lyfter sin hand.\\
Ännu Svearnes mod sif bereder en Stod\\
Utaf Lagrarne höjd, fast besglad mes blod.\\
Himlen gifve oss frid!\\
Men om den icke vinne utan vapen och strid,\\
Blifve Segern då vår, ögat offre sin tår\\
Den som faller, ock rätta till vår tacksamhet får.\\
\\
{\footnotesize\textit{Text: Samuel Ödmann\\
Musik: Johann Christian Friedrich Haeffner\\
\\
Anses vara den första upsaloensiska studentsången. Musiken hade
skrivits tidigare som soldatkör till Hæffners opera “Renaud” (\oldstylenums{1801}).\\
\\
Framfördes första gången, mitt under kriget mot Ryssland,
den \oldstylenums24} oktober \oldstylenums{1808}. Detta skedde vid en
uppvakting för fältmarskalkten Klingspor på genomresa från Finland
till Stockholm som gästade landshövding Uppsala slott.  Uppsalas
studenter marscherade då upp till slottet. Sångerna gick i täten och
sjöng denna sång.}}


\chapterpage{Datavetarvisor}{Datavetarvisor}
\newpage
\thispagestyle{empty}
\noindent{\texttt{Datavetare är skapande människor\\
\\
Datavetare skapar. De skapar nya program som världen aldrig sett
tidigare, för att låta datorer göra saker som de aldrig gjort
tidigare.  De forskar fram nya kunskaper och öppnar nya fält för
framtidens forskning. Datavetare skapar de verktyg som gör det möjligt
att lösa nya problem.  Nya programmeringsspråk, nya beräkningsverktyg,
bättre kompilatorer, bättre databaser, nya operativsystem, bättre
utvecklingsverktyg, distribuerade system... Listan kan göras hur lång
som helst.\\
\\
Datavetenskapens handfasta mål är att förbättra den programvara som vi
kör på våra datorer. Vad som krävs är att programmen gör vad de skall
(korrekthet), att de inte gör fel om något oförutsett inträffar
(pålitlighet), och att de arbetar effektivt (effektivitet). Dessutom
skall programmen utvecklas på ett effektivt sett.}}\\
\\
{\footnotesize\textit{Ur ``Vad är en datavetare egentligen''\\
    BlurGel nr 0, 1995}}
\clearpage
\svisa{THE BASIC SONG@\texttt{THE BASIC SONG}}{10 LET oss nu fatta i våra glas@\texttt{10 LET oss nu fatta i våra glas}}{\texttt{THE BASIC SONG}}
{\footnotesize\textit{Melodi: Mors lilla Olle}}\\
\\
\texttt{10 LET oss nu fatta i våra glas}\\
\texttt{20 INPUT en klunk utav det som där has}\\
\texttt{30 IF du fått nog THEN till 50 min vän}\\
\texttt{40 ELSE GOTO-baka till 10 igen}\\
\texttt{50 END}\\
\\
{\footnotesize\textit{Inledande radnummer sjungs ej.\\ I övrigt följs
    kommandon, ex. drick vid ``\texttt{INPUT en klunk}...''}}


\chapterpage{Snapsvisor}{Snapsvisor}

\chapterpage{Ölvisor}{Ölvisor}
%\customindex{YOLO}{49}
\chapterpage{Vinvisor}{Vinvisor}

\chapterpage{Avecvisor}{Avecvisor}

\chapterpage{Gasquevisor}{Gasquevisor}

\chapterpage{Säsongsvisor}{Säsongsvisor}

\chapterpage{Svenska visor}{Svenska visor}

\chapterpage{Osvenska visor}{Osvenska visor}

\chapterpage{Phulavisor}{Phulavisor}

\chapterpage{Register}{Register}

\cfoot{}
\newgeometry{top=1cm,bottom=1.2cm,left=1cm,right=1cm,footskip=0pt,headsep=0.2cm}
\small
\printindex
\hiddenchapterpage{\texttt{\emptyset}}
\ohead{}
\ofoot[\pagemark]{}
\newpage
\null
\newpage
\null
\newpage
\null
\newpage
\null
\newpage
\null
\newpage
\null
\newpage
\null
\newpage
\null
\end{document}
